\documentclass[12pt]{article}
% Préambule commun à tous les documents extraits de « La Revue CAPES »
\usepackage[T1]{fontenc}
\usepackage[utf8]{inputenc}
\usepackage{lmodern}
\usepackage{amsmath,amssymb,amsthm,amsfonts,mathrsfs}
\usepackage{setspace}
\usepackage{listings}
\usepackage{pgfplots}
\pgfplotsset{compat=1.18}
\usepackage{longtable}
\usepackage{adjustbox}
\usepackage{lettrine}
\usepackage{tkz-tab}
\usepackage{changepage}
\usepackage{pdflscape}
\usepackage{graphicx}
\usepackage{tikz}
\usepackage{xcolor}
\usepackage[a4paper,top=2cm,bottom=2.5cm,left=2.5cm,right=2cm]{geometry}
\usepackage{fancyhdr}
\usepackage{draftwatermark}
\SetWatermarkText{La RevueCapes}
\SetWatermarkLightness{0.9}
\SetWatermarkScale{2}
\usepackage{hyperref}
\hypersetup{colorlinks=true,linkcolor=blue!50!black,urlcolor=blue!50!black}

\definecolor{revue}{RGB}{20,60,120}

\pagestyle{fancy}
\fancyhf{}
\renewcommand{\headrulewidth}{0.4pt}
\setlength{\headheight}{15pt}

% \entete{Matière}{Titre}{Type de document}
\newcommand{\entete}[3]{%
  \fancyhead[L]{\small\textsc{La Revue CAPES} -- #1}%
  \fancyhead[R]{\small #2}%
  \fancyfoot[C]{\thepage}%
  \thispagestyle{plain}%
  \begin{center}
    {\color{revue}\rule{\textwidth}{1.2pt}}\\[4pt]
    {\small\textsc{Concours directs CAP-CEG / CAPES -- Burkina Faso}}\\[6pt]
    {\LARGE\bfseries\color{revue} #2}\\[6pt]
    {\large\itshape #1 -- #3}\\[2pt]
    {\color{revue}\rule{\textwidth}{1.2pt}}
  \end{center}
  \vspace{0.5cm}%
}

% Encadré pour les remarques ajoutées dans les corrigés
\newcommand{\remarque}[1]{\par\medskip\noindent\fcolorbox{revue}{revue!6}{\parbox{\dimexpr\linewidth-2\fboxsep-2\fboxrule}{\textbf{\color{revue}Remarque.} #1}}\par\medskip}


\begin{document}
\entete{Mathématiques}{Corrigé détaillé -- Session 2021}{Correction}
\begin{spacing}{1.5}

\textbf{Exercice 1}

Étudions la nature(convergente ou divergente) des séries numériques suivantes : 

\[\bullet S_1=\sum^{+\infty}_{n=1}(1-\frac{1}{n})^n\]

Soit $u_n=(1-\frac{1}{n})^n\Longrightarrow S_1=\sum^{+\infty}_{n=1}u_n$

Calculons la limite de $u_n$

On a :

$u_n=(1-\frac{1}{n})^n=e^{nln(1-\frac{1}{n})}$

Si $n\longrightarrow +\infty$ alors $-\frac{1}{n}\longrightarrow 0$ et $ln(1-\frac{1}{n})\sim -\frac{1}{n}+\circ(\frac{1}{n})$ (Développement limité de Taylor)


Ainsi,

 \[\lim_{n\to +\infty} u_n=\lim_{n\to +\infty}e^{n(-\frac{1}{n}+\circ(\frac{1}{n}))}=\lim_{n\to +\infty}e^{-1+\circ(1)}\longrightarrow \frac{1}{e}\neq 0\] alors $S_1$ diverge.

\medskip
\[\bullet S_2=\sum^{+\infty}_{n=1}\frac{n(n-1)}{(n+1)!}\].

Soit $u_n=\frac{n(n-1)}{(n+1)!}\Rightarrow u_{n+1}=\frac{n(n+1)}{(n+2)!}$

On a : 
\begin{align*}
\lim_{n\to +\infty}\frac{u_{n+1}}{u_n} &=\lim_{n\to +\infty}\frac{\frac{n(n+1)}{(n+2)!}}{\frac{n(n-1)}{(n+1)!}}\\
&=\lim_{n\to +\infty}\frac{n(n+1)(n+1)!}{n(n-1)(n+2)!}\\
&=\lim_{n\to +\infty}\frac{n(n+1)(n+1)!}{n(n-1)(n+2)(n+1)!}\\
&=\lim_{n\to +\infty}\frac{n+1}{(n-1)(n+2)}=\lim_{n\to +\infty}\frac{n}{n^2}=\lim_{n\to +\infty}\frac{1}{n}=0<1
\end{align*}, alors $S_2$ converge d'après le critère d'Alembert.



\medskip
\textbf{Exercice 2 }



1 : Développement en série entière de la fonction \( u \)

On considère l'équation différentielle suivante :
\[
(x - 1) u'(x) + u(x) = 0,
\]
avec la condition initiale \( u(0) = -1 \). Nous cherchons le développement en série entière de la solution \( u(x) \).

\quad
 Résolution de l'équation différentielle

Nous réécrivons l'équation sous la forme :
\[
u'(x) = \frac{-u(x)}{x - 1}.
\]

Nous cherchons une solution sous la forme d'une série entière autour de \( x = 0 \), soit :
\[
u(x) = \sum_{n=0}^{\infty} a_n x^n.
\]

En dérivant cette série terme à terme, on obtient :
\[
u'(x) = \sum_{n=1}^{\infty} a_n n x^{n-1}.
\]

Nous remplaçons cette expression dans l'équation différentielle :
\[
(x - 1) \sum_{n=1}^{\infty} a_n n x^{n-1} + \sum_{n=0}^{\infty} a_n x^n = 0.
\]

\quad
 Développement des termes

Développons \( (x - 1) \) et multiplions la série par \( (x - 1) \). Nous obtenons :
\[
(x - 1) \sum_{n=1}^{\infty} a_n n x^{n-1} = \sum_{n=1}^{\infty} a_n n (x^{n} - x^{n-1}).
\]

Cela donne :
\[
\sum_{n=1}^{\infty} a_n n x^n - \sum_{n=1}^{\infty} a_n n x^{n-1}.
\]

En ajoutant la deuxième série \( \sum_{n=0}^{\infty} a_n x^n \), on obtient :
\[
\left( \sum_{n=1}^{\infty} a_n n x^n - \sum_{n=1}^{\infty} a_n n x^{n-1} \right) + \sum_{n=0}^{\infty} a_n x^n = 0.
\]

Nous réorganisons les indices et combinons les termes en \( x^n \), ce qui nous permet d'obtenir une relation de récurrence entre les coefficients \( a_n \).

\quad
 Résolution de la relation de récurrence

Nous trouvons la relation de récurrence suivante pour les coefficients \( a_n \) :
\[
a_n = \frac{1}{n} a_{n-1}, \quad \text{pour } n \geq 1.
\]

Nous utilisons la condition initiale \( u(0) = -1 \), ce qui donne \( a_0 = -1 \). En utilisant la relation de récurrence, nous pouvons calculer les autres coefficients :
\[
a_1 = \frac{1}{1} a_0 = -1, \quad a_2 = \frac{1}{2} a_1 = \frac{1}{2}, \quad a_3 = \frac{1}{3} a_2 = \frac{1}{6}, \quad \dots
\]

Ainsi, la série entière de la solution \( u(x) \) est :
\[
u(x) = \sum_{n=0}^{\infty} \frac{(-1)^n}{n!} x^n.
\]

Il s'agit du développement en série entière de la fonction \( u(x) \).

 2. Rayon de convergence de la série

Le rayon de convergence \( R \) de la série entière \( u(x) = \sum_{n=0}^{\infty} a_n x^n \) est donné par la formule de Cauchy-Hadamard :
\[
\frac{1}{R} = \limsup_{n \to \infty} \sqrt[n]{|a_n|}.
\]

Ici, nous avons \( a_n = \frac{(-1)^n}{n!} \), donc :
\[
\limsup_{n \to \infty} \sqrt[n]{|a_n|} = \limsup_{n \to \infty} \sqrt[n]{\frac{1}{n!}} = 0.
\]

Ainsi, le rayon de convergence est \( R = \infty \).

 3. Domaine de convergence

Étant donné que le rayon de convergence de la série est \( R = \infty \), cela signifie que la série converge pour tous les \( x \in \mathbb{R} \). Ainsi, le domaine de convergence de la série est \( \mathbb{R} \).



\medskip
\textbf{Exercice 3}

Calculons $\int\int_{D}f(x,y)dxdy$ dans les cas suivants :

1. $D= \lbrace {y\geq 0, x+y, y-x\leq 1\rbrace}$ et $f(x,y)=x^{2}y$.


On a:

$
\begin{cases}
x+y\leq 1 (1)\\
y-x \leq 1 (2)
\end{cases}$

$(1)+(2)\Longrightarrow 2y\leq 2 \Longrightarrow y\leq 1$

De plus, $y\geq 0$. 	D'où $0\leq y\leq 1\Rightarrow y\in [0,1]$

$
\begin{cases}
x+y\leq 1 \\
y-x \leq 1 
\end{cases}\Rightarrow 
\begin{cases}
x\leq 1-y \\
x \geq y-1
\end{cases}\Longrightarrow  y-1\leq x \leq 1-y \Rightarrow x\in [y-1,1-y]$

Donc,
$
\int\int_{D}f(x,y)dxdy=\int_0^1\int_{y-1}^{1-y}x^2ydxdy=\int_0^1(\int_{y-1}^{1-y}x^2y.dx)dy
$

$\int_{y-1}^{1-y}x^2y.dx=2\int_{0}^{1-y}x^2y.dx$ car $x\mapsto x^2y$ est une fonction paire pour $y$ fixée.

Donc

\begin{align*}
\int\int_{D}f(x,y)dxdy &=\int_0^1(2\int_{0}^{1-y}x^2ydx)dy\\
&=2\int_0^1([\frac{1}{3}x^3y]_0^{1-y})dy\\
&=\frac{2}{3}\int_0^1y(1-y)^3dy\\
&=\frac{2}{3}\int_0^1y(1-3y+3y^2-y^3)dy\\
&=\frac{2}{3}\int_0^1(y-3y^2+3y^3-y^4)dy\\
&=\frac{2}{3}[\frac{1}{2}y-y^3+\frac{3}{4}y^4-\frac{1}{5}y^5]_0^1\\
&=\frac{2}{3}(\frac{1}{2}-1+\frac{3}{4}-\frac{1}{5})\\
\int\int_{D}f(x,y)dxdy &=\frac{1}{30}
\end{align*}

\medskip
2. $D= \lbrace 0\leq y\leq x, x^2+y^2 \leq 1\rbrace$ et $f(x,y)=e^{x^2+y^2}$.

Posons :

$x=rcos\theta$ et $y=rsin\theta$

$x^2+y^2=r^2$

$x^2+y^2 \leq 1\Rightarrow r^2\leq 1\Longrightarrow 0\leq r \leq 1\Longrightarrow r\in [0,1]$

$0\leq y\leq x\Rightarrow 0\leq rsin\theta\leq rcos\theta\Rightarrow 0\leq sin\theta\leq cos\theta\\Rightarrow 0\leq \tan\theta\leq 1$. En appliquant $\arctan(x)$ à cette inégalité, on obtient :

$\arctan(0)\leq \arctan(\tan(\theta))\leq \arctan(1)\Longrightarrow 0\leq \theta\leq \frac{\pi}{4}\Longrightarrow \theta\in [0,\frac{\pi}{4}]$

Calculons le jacobien :

$j(r,\theta)=\begin{vmatrix}
\frac{\partial x}{\partial r} & \frac{\partial x}{\partial \theta}\\
\frac{\partial y}{\partial r} & \frac{\partial x}{\partial \theta}
\end{vmatrix}=\begin{vmatrix}
cos \theta & -r\sin\theta\\
\sin\theta & r\cos\theta
\end{vmatrix}=r\Rightarrow dxdy=rdrd\theta$

Donc,

\begin{align*}
\int\int_{D}f(x,y)dxdy &=\int_0^{\frac{\pi}{4}}\int_{0}^1e^{r^2}rdrd\theta\\
&=\int_0^{\frac{\pi}{4}}d\theta\times\int_{0}^1re^{r^2}dr\\
&=[\theta]_0^{\frac{\pi}{4}}\times[\frac{1}{2}e^{r^2}]_0^1\\
&=\frac{\pi}{4}\times\frac{1}{2}(e-1)=\frac{\pi}{8}(e-1)
\end{align*}

\medskip
\textbf{Exercice 4}

Soit l'équation différentielle

$(E) : y^{\prime\prime}-4y^{\prime} + 4y = (t^2 - 1)e^t$.

1. Trouvons une solution particulière de $(E)$ sous la forme $y_p= Q(t)e^t$ ou $Q$ est un polynôme de degré 2.

$y_p=(at^2+bt+c)e^t$

$y'_p=(2at+b)e^t+(at^2+bt+c)e^t=(at^2+(2a+b)t+b+c)e^t$

$y''_p=(2at+2a+b)e^t+(at^2+(2a+b)t+b+c)e^t=(at^2+(4a+b)t+2a+2b+c)e^t$

En reportant dans l'équation différentielle $(E)$, on obtient;

$(at^2+(4a+b)t+2a+2b+c)e^t-4(at^2+(2a+b)t+b+c)e^t+4(at^2+bt+c)e^t=(t^2 - 1)e^t$

$(at^2+(b-4a)t+(2a-2b+c))e^t=(t^2 - 1)e^t$

Par identification,

$\begin{cases}
a=1\\
b-4a=0\\
2a-2b+c=-1
\end{cases}\Rightarrow \begin{cases}
a=1\\
b=4\\
c=5
\end{cases}$

D'où $y_p=(t^2+4t+5)e^t$

2. Déduisons la solution générale de $(E)$.

Calculons la solution homogène :

L'équation homogène est donnée par :

$y^{\prime\prime}-4y^{\prime} + 4y =0$

L'équation caractéristique :

$r^2-4r+4=0\Rightarrow (r-2)^2=0\Rightarrow r_0=2$

On a une racine double, donc la solution homogène est de la forme :

$y_0=(At+B)e^{2t}$ avec $A,B\in \mathbb{R}$

Ainsi la solution générale est donnée par :

\[
y_G=y_0+y_p=(At+B)e^{2t}+(t^2+4t+5)e^t\] avec $A,B\in \mathbb{R}$


3. Trouvons la solution $h$ telle que $h(0)=h^{\prime}(0)=1$

$h(t)=(At+B)e^{2t}+(t^2+4t+5)e^t$

$h(0)=1\Rightarrow B+5=1\Rightarrow B=-4$

$h'(t)=Ae^{2t}+2(At+B)e^{2t}+(2t+4)e^t+(t^2+4t+5)e^t$

$h'(0)=1\Rightarrow A+2B+4+5=1\Rightarrow A+2B=-8\Rightarrow A=0$

$h(t)=-4e^{2t}+(t^2+4t+5)e^t$

\medskip

\end{spacing}
\end{document}
